\section{Exact messages at fixed treewidth}\label{sec:spectral}
Section~\ref{sec:main-result} outlined the proof of
Theorem~\ref{thm:main}. This section completes it by meeting the two
requirements of Theorem~\ref{thm:dictionary} for indicator quadratics: a
Lipschitz constant common to all branches and a restricted additive
oracle. Section~\ref{sec:spectral-bounds} bounds all conditional
optimizers uniformly, which gives the Lipschitz constant, and
Section~\ref{sec:spectral-oracle} uses these bounds to build the oracle.
Section~\ref{sec:spectral-main} proves Theorem~\ref{thm:main}, and
Section~\ref{sec:spectral-certificate} bounds the original optimum.

\subsection{Uniform bounds on conditional optimizers}
\label{sec:spectral-bounds}
Both requirements need bounds that hold for all supports at once. The
net of Theorem~\ref{thm:dictionary} has radius $\varepsilon/(2L)$, so $L$
must be a Lipschitz constant for every branch. The oracle will round a
restricted optimizer to a finite grid, and the rounded point is close to
the optimizer only if the grid covers a range that is known in advance
to contain it. The oracle does not know which support is optimal, so
this range must contain the conditional optimizers of all supports.

Let $k=|S|$ and $A\subseteq I$. Principal submatrices inherit
\eqref{eq:spectral}, so $\|Q_{AA}^{-1}\|_2\le1/\mu$. Moreover
$\|Q_{AS}\|_2\le\|Q\|_2\le H$, $\|c_A\|_2\le C\sqrt n$, and
$\|t\|_2\le\sqrt kR$ on $[-R,R]^k$. Inserting these bounds into
\eqref{eq:optimizer}, and into the gradient of \eqref{eq:branch}, shows
that every $t\in[-R,R]^k$ satisfies
\begin{align}
 \|x_A(t)\|_2&\le\frac{C\sqrt n+2H\sqrt{k}R}{2\mu}
                    \le X_k:=\frac{nC+2HkR}{2\mu},\label{eq:boundx}\\
 \|\nabla q_A(t)\|_2&=\|2Q_{SA}x_A(t)\|_2
                    \le L_k:=2HX_k.\label{eq:boundL}
\end{align}
The second bound in \eqref{eq:boundx} replaces the square roots by a
rational quantity, which can be computed exactly. Since the box is
convex, \eqref{eq:boundL} makes every $q_A$, and hence every $F_A$,
$L_k$-Lipschitz on it. The bounds are uniform over supports, so they
still hold when some indicator bits are fixed: a restricted optimum is
the conditional optimizer of some support.

The bounds grow with $R$ because the boundary values enter the linear
term $c_A+2Q_{AS}t$ of the conditional problem. For the same reason, a
bound on the global optimizer alone would not suffice: it need not bound
the conditional optimizers for a prescribed boundary value.

The same estimate also shows how large the box must be to contain every
optimal boundary value. For every $\xi$, each optimal
solution of \eqref{eq:problem} is the conditional optimizer of its
support for $I=[n]$ and $S=\varnothing$, so \eqref{eq:boundx} with $k=0$
gives $\|x\|_2\le C\sqrt n/(2\mu)\le nC/(2\mu)$. Hence, for every
separator $S$, the box with rational radius $R=nC/(2\mu)$ contains the
boundary values $x_S$ of all optimal solutions.

\subsection{A restricted grid oracle}\label{sec:spectral-oracle}
Fix a message $M_{I,S}$, a rational boundary point $t\in[-R,R]^k$, and
an accuracy $\varepsilon>0$. The oracle must approximately minimize the
internal objective of \eqref{eq:message} at $t$ under arbitrary fixed
indicator bits, as \eqref{eq:oracle} requires. If the continuous
variables could take only finitely many values, this would be a discrete
problem whose costs each involve one vertex or one edge of the support
graph, and dynamic programming over the tree decomposition would solve
it exactly. We therefore restrict every coordinate to a finite grid,
fine enough that rounding costs at most $\varepsilon$, solve the grid
problem exactly, and return the support of the grid solution with its
exact continuous cost.

First suppose $X_k>0$. Let $K$ be the least positive even integer
satisfying
\begin{equation}\label{eq:gridB}
 K^2\varepsilon\ge HnX_k^2.
\end{equation}
Then $K\le2+X_k\sqrt{Hn/\varepsilon}$. Use the coordinate grid
\[
 G_K=\{-X_k+2X_kj/K:0\le j\le K\}.
\]
Its $K+1$ points are spaced $2X_k/K$ apart, so every point of
$[-X_k,X_k]$ lies within $X_k/K$ of the grid. The proof of
Lemma~\ref{lem:gridoracle} shows that \eqref{eq:gridB} then makes the
cost of rounding at most $\varepsilon$. Each vertex $i\in I$ has an
inactive state $(z_i,x_i)=(0,0)$ and active states $(1,a)$ for
$a\in G_K$. Since $K$ is even, $0\in G_K$: as in the continuous problem,
an active coordinate may be zero, and this state is distinct from the
inactive state. Fixed bits delete incompatible states. For example,
$K=4$ gives $G_4=\{-X_k,-X_k/2,0,X_k/2,X_k\}$ and six states per vertex;
fixing $z_i=1$ leaves the five active states, and fixing $z_i=0$ leaves
only the inactive one.

The grid objective is the sum of the unary costs
\[
 Q_{ii}x_i^2+(c_i+2(Q_{IS}t)_i)x_i+(\lambda_i+\xi_i)z_i
\]
and the edge costs $2Q_{ij}x_ix_j$ for $i<j$ in $I$ with $Q_{ij}\ne0$,
evaluated at grid states. These are the terms of the internal objective
of \eqref{eq:message} at $t$; the boundary enters only through the linear
coefficients.

The dynamic program works on the rooted decomposition of
Section~\ref{sec:prelim-treedec}. Restrict every bag to $I$, and write
$B_u^I=B_u\cap I$ for the restricted bag of node $u$. Restricting the
bags to $I$ preserves the defining conditions of a tree decomposition,
so the restricted bags form a tree decomposition of the subgraph induced
by $I$; empty bags are harmless. Assign each unary and edge cost to
exactly one bag containing its variables. For a bag $u$ and an
assignment $s$ of states to its vertices, let $a_u(s)$ be the sum of its
assigned costs, and let $D_u(s)$ be the minimum total of the costs
assigned to the subtree of $u$, over assignments that agree with $s$ on
$B_u^I$. The recurrence is
\begin{equation}\label{eq:DP}
 D_u(s)=a_u(s)+\sum_{v\text{ child of }u}
       \min_{r:\ r|_{B_u^I\cap B_v^I}=s|_{B_u^I\cap B_v^I}}D_v(r).
\end{equation}
The inner minimum is a discrete analogue of a message: for each
assignment of states to the shared vertices $B_u^I\cap B_v^I$, it gives
the least total of the costs assigned to the subtree of $v$. Because each
vertex has finitely many states, it is a finite table rather than a
minimum of quadratics in continuous variables. The running-intersection
property makes the internal sets of different children disjoint, and
it ensures that agreement on $B_u^I\cap B_v^I$ is sufficient for
compatibility. Induction from the leaves shows that \eqref{eq:DP}
computes $D_u$, and minimizing $D$ at the root gives the exact grid
optimum.

Each child table is first minimized over its states at each assignment
to the shared vertices $B_u^I\cap B_v^I$; reading these projected tables
at the parent avoids a product of two bag tables. A restricted bag has at
most $w+1$ vertices with at most $K+2$ states each, so each table has at
most $(K+2)^{w+1}$ entries. Each entry adds its assigned costs and one
projected value per child, and each child table is projected once. With
$b$ bags and $O(n^2)$ costs, the number of arithmetic operations is
\begin{equation}\label{eq:DPcost}
 O\bigl((b+n^2)(K+2)^{w+1}\bigr).
\end{equation}
Back-pointers recover a minimizing grid solution and its support. The
bound needs no balance or degree condition on the decomposition.

\begin{lemma}[Restricted additive oracle]\label{lem:gridoracle}
Under any fixed-bit restriction, solve the grid problem and let $A$ be the
support of its optimal solution. Returning $A$ with its exact perturbed
value $F_A(t)$ from \eqref{eq:noisybranch} satisfies \eqref{eq:oracle}.
\end{lemma}

\emph{Proof idea.} Rounding a true restricted optimizer to the grid raises
its cost by at most $\varepsilon$, because the linear part of the
rounding error vanishes at a minimizer. The grid optimum is at least as
good, and reoptimizing its support continuously can only lower its cost.

\begin{proof}
Let $A^*$ be an optimal support of the restricted continuous problem and
$x^*$ its unique continuous minimizer, which is $x_{A^*}(t)$ extended by
zeros. By \eqref{eq:boundx}, $x^*\in[-X_k,X_k]^I$. Round each active
coordinate to a nearest point of $G_K$ and keep the inactive coordinates
zero. The rounding error $e$ is supported on $A^*$ and has
$|e_i|\le X_k/K$, so $\|e\|_2^2\le nX_k^2/K^2$. Since $x^*$ is
stationary on $A^*$, the linear term in $e$ cancels, and the objective
increases by
\[
 e^TQ_{II}e\le H\|e\|_2^2\le HnX_k^2/K^2\le\varepsilon
\]
by \eqref{eq:gridB}. The rounded point has support $A^*$, so it obeys the
fixed bits, and the optimal grid solution costs no more. Its support $A$
also obeys the fixed bits, and continuous reoptimization on $A$ can only
decrease the cost, so $F_A(t)$ is at most the grid optimum. Hence $F_A(t)$
lies between the restricted optimum and the restricted optimum plus
$\varepsilon$.
\end{proof}

If $X_k=0$, then \eqref{eq:boundx} forces every conditional optimizer to
be zero, so $F_A(t)=\sum_{i\in A}(\lambda_i+\xi_i)$. An exact oracle then
keeps the fixed bits and sets each free bit to minimize its own perturbed
penalty $\lambda_i+\xi_i$.

The construction and the lemma use $X_k$ only to ensure that
$x^*\in[-X_k,X_k]^I$. They therefore remain valid when $X_k$ is replaced
throughout by any larger rational bound, which lets
Theorem~\ref{thm:main} use one common bound for all messages. Every
fixed-bit restriction is feasible, because an active coordinate may be
zero, so the oracle never reports an empty cell.

\subsection{The algorithm and its expected bit complexity}
\label{sec:spectral-main}
Let the supplied decomposition have width $w$ and $b$ bags, and let the
box radius $R\ge0$ and the perturbation scale $\sigma>0$ be rational.
The algorithm uses one perturbation draw and one set of parameters for
all messages, chosen so that Theorem~\ref{thm:dictionary} and
Lemma~\ref{lem:gridoracle} apply to every message at once. Every
separator has $|S|\le w+1$, and $X_k$ and $L_k$ do not decrease with $k$,
so it suffices to evaluate \eqref{eq:boundx} and \eqref{eq:boundL} at
$k=w+1$. With $\phi=1/(2\sigma)$ from \eqref{eq:gridmass}, the accuracy
\eqref{eq:accuracy} becomes $\sigma/(6n)$. This gives the parameters
\begin{equation}\label{eq:mainparameters}
\begin{gathered}
 X=\frac{nC+2H(w+1)R}{2\mu},\quad L=2HX,\quad
 \varepsilon=\frac{\sigma}{6n},\quad
 J_*=\Bigl(1+\frac{12n(w+1)RL}{\sigma}\Bigr)^{w+1},\\
 K=\text{the least positive even integer with }K^2\varepsilon\ge HnX^2.
\end{gathered}
\end{equation}
Thus $K$ is defined by \eqref{eq:gridB} with $X$ in place of $X_k$, and
$K=2$ if $X=0$. By \eqref{eq:netsize} with $\varepsilon=\sigma/(6n)$ and
$k\le w+1$, $J_*$ bounds the number of points of the midpoint net on
every box.

For each message, the algorithm runs certified enumeration
(Lemma~\ref{lem:enumeration}) with tolerance $\varepsilon$ at every
point of the midpoint parameter net of the box
(Section~\ref{sec:net-work}), using the oracle of
Lemma~\ref{lem:gridoracle} with $X$ in place of $X_k$. It keeps the
union of the supports found and builds the perturbed branch $F_A$ of
each. Every oracle call runs the grid dynamic program on all of $I$.

\begin{theorem}[Exact messages at fixed treewidth]\label{thm:main}
Fix $w$. Assume \eqref{eq:problem}--\eqref{eq:spectral}, a supplied tree
decomposition of width $w$ with $b$ bags, and rational $R\ge0$ and
$\sigma>0$. Let $N$ be the smallest power of two with $N\ge2n$, and let
$\xi_1,\ldots,\xi_n$ be independent and uniform on the perturbation grid
\eqref{eq:noisegrid}. There is an algorithm with the following
properties.
\begin{enumerate}
\item For every realization of $\xi$, it returns, for every message
 $M_{I,S}$ of Section~\ref{sec:prelim-messages}, a dictionary of
 quadratics with rational coefficients on the box $[-R,R]^{|S|}$. Each
 dictionary contains every support that is optimal somewhere on its box.
 The algorithm also returns a global optimizer of the perturbed problem
 \eqref{eq:problem}.
\item It uses $n\log_2N=O(n\log(n+1))$ random bits.
\item Let $\ell$ be the total binary input length, and let $X$, $J_*$,
 and $K$ be the parameters \eqref{eq:mainparameters}. The expected
 number of bit operations of the algorithm is at most
 \begin{equation}\label{eq:explicitwork}
  P_w(\ell)\,bJ_*\,(K+2)^{w+1},\qquad
  K+2\le4+Xn\sqrt{6H/\sigma},
 \end{equation}
 where $P_w$ is a polynomial that depends only on $w$.
\end{enumerate}
In particular, at fixed $w$ the expected bit complexity is polynomial in
$\ell$ and in the numerical quantities $C,H,\mu^{-1},R,\sigma^{-1}$.
\end{theorem}

The factors in \eqref{eq:explicitwork} follow the structure of the
algorithm: at most $b$ messages, at most $J_*$ net points per message,
an expected $O(n)$ oracle calls per net point (absorbed into
$P_w(\ell)$), and $(K+2)^{w+1}$ states per bag in each call. Each
dictionary is the union of the enumeration outputs at the net points of
its box, so by \eqref{eq:dictionarysize} its expected size is at most
$eJ_*$.

\emph{Proof idea.} Correctness for every realization of $\xi$ follows
from Theorem~\ref{thm:dictionary} and Lemma~\ref{lem:gridoracle}. The
rest of the proof checks that all numbers have polynomial bit length, so
that operation counts become bit counts, and sums the expected oracle
calls over the messages.

\begin{proof}
\emph{Step 1: correctness.} Fix a message $M_{I,S}$, so $k=|S|\le w+1$.
Apply Theorem~\ref{thm:dictionary} with $Z=\{0,1\}^I$, $m=|I|\le n$,
$T=[-R,R]^k$, the deterministic functions $q_A$ of \eqref{eq:branch}
(so that the perturbed costs are the branches $F_A$),
$\phi=1/(2\sigma)$, and $\tau=1/N$. These values satisfy
\eqref{eq:intervalmass} by \eqref{eq:gridmass}, the choice of $N$ gives
$\tau\le1/(2n)$, and \eqref{eq:accuracy} gives $\varepsilon=\sigma/(6n)$.
Since $k\le w+1$, the common bounds $X$ and $L$ dominate
\eqref{eq:boundx} and \eqref{eq:boundL}, so every $q_A$ is $L$-Lipschitz
on the box. Use the midpoint net of \eqref{eq:netsize}, or one point if
the box is a singleton or $L=0$. Its cardinality is at most
$(1+12nkRL/\sigma)^k\le J_*$. At each net point, the oracle is
Lemma~\ref{lem:gridoracle} with $X$ in place of $X_k$, or the exact
oracle for the case $X=0$. Ties in the dynamic program and in the
enumeration are broken by fixed orders. Since $\xi$ takes finitely many
values, the measurability required by Theorem~\ref{thm:dictionary} holds
automatically. Let $\mathcal D$ be the union of the outputs. By
\eqref{eq:dictionaryexact} and \eqref{eq:envelope}, the branches $F_A$,
$A\in\mathcal D$, form a dictionary for $M_{I,S}$ on the box that contains
every support optimal somewhere on it. This holds for every realization.
For $M_{[n],\varnothing}$ the branches are constants. The smallest one
gives a global minimizing support, and \eqref{eq:optimizer} gives its
rational optimizer.

\emph{Step 2: bit lengths of parameters and dynamic-program values.}
Every input rational has bit length at most $\ell$, and $n,b\le O(\ell)$.
The quantities $X$, $L$, $\varepsilon$, the subdivision count of the net,
the net coordinates, and the sampled penalties are formed by sums,
products, reciprocals of positive rationals, and integer ceilings. Hence
they have polynomial bit length. Binary search with exact rational
comparisons finds $K$ in polynomially many bit operations, and $\log K$
is polynomial; no square root is computed. Every grid point is an
integer multiple of $X/K$ with absolute value at most $X$. Hence each
unary or edge cost takes all its grid values with one fixed denominator
and a bounded magnitude, both of polynomial bit length. A finite entry of
\eqref{eq:DP} is a sum of values of a subset of these $O(n^2)$ costs,
each used at most once. The product of the cost denominators is a common
denominator, and both it and the entry's magnitude have polynomial bit
length. Minima select values and create no new denominators. Hence all
additions and comparisons in the dynamic program have polynomial bit
cost, for every realization of $\xi$, including realizations with large
enumerations.

\emph{Step 3: branch construction.} Each exact value $F_A(t)$ returned by
the oracle, and each branch $F_A$, comes from a principal system
$Q_{AA}u=d$. Clearing denominators gives an integer system whose entries
have polynomial bit length. Each determinant is a sum of at most $n!$
products of at most $n$ such integers, so its bit length is polynomial
in $\ell$. Since $Q_{AA}$ is nonsingular, exact rational elimination
gives the coefficients of \eqref{eq:optimizer} and \eqref{eq:branch} with
polynomial bit length in polynomial bit time. Adding the rational
perturbations gives the same bound for $F_A$. The affine map
\eqref{eq:optimizer} can be stored with each branch to recover
conditional optimizers.

\emph{Step 4: expected work.} By \eqref{eq:dictionarysize} with
$J\le J_*$ and $m\le n$, the expected number of oracle calls is at most
$J_*(1+en)$ per message, hence at most $bJ_*(1+en)$ in total. Linearity
of expectation applies even though the messages share $\xi$ and their
subproblems overlap. Each call uses \eqref{eq:DPcost} arithmetic
operations, each of polynomial bit cost by Step~2. Heap operations,
deduplication, exact support evaluations, and branch construction add a
polynomial factor per output or call (Section~\ref{sec:net-work}), and
each output comes from a call. The factors $b+n^2$ and $1+en$ and these
polynomial factors are absorbed into $P_w(\ell)$. Finally,
\eqref{eq:gridB} gives $K\le2+X\sqrt{Hn/\varepsilon}$, and
$\varepsilon=\sigma/(6n)$ turns this into the bound on $K+2$ in
\eqref{eq:explicitwork}.
\end{proof}

The dictionaries are exact on the box and need not be exact outside it,
where the least retained branch may exceed the message. The randomness
affects only the running time. For exceptional realizations of $\xi$ the
dictionaries may be exponentially large, but the algorithm still
terminates with exact output.

\subsection{An additive certificate for the original objective}
\label{sec:spectral-certificate}
Theorem~\ref{thm:main} is exact for the perturbed penalties
$\lambda+\xi$, but the original problem has penalties $\lambda$. Let
$v_0$ be its optimum, the value of \eqref{eq:problem} at $\xi=0$. The
perturbation changes the cost of each solution $(x,z)$ by $\xi^Tz$, and
$|\xi_i|\le\sigma$, so an exact perturbed optimizer is also nearly
optimal for the original problem. Let $(\widehat x,\widehat z)$ be any
exact optimizer of the perturbed problem, however computed, and define
\begin{equation}\label{eq:certificate}
 \mathrm{LB}=v_\xi-\sum_i\max\{\xi_i,0\},\qquad
 \mathrm{UB}=v_\xi-\xi^T\widehat z.
\end{equation}
For every feasible $(x,z)$, the unperturbed value equals the perturbed
value minus $\xi^Tz$. The perturbed value is at least $v_\xi$ and
$\xi^Tz\le\sum_i\max\{\xi_i,0\}$, so the unperturbed value is at least
$\mathrm{LB}$. At $(\widehat x,\widehat z)$ the unperturbed value equals
$\mathrm{UB}$. Hence
\begin{equation}\label{eq:certgap}
 \mathrm{LB}\le v_0\le\mathrm{UB},\qquad
 \mathrm{UB}-\mathrm{LB}
 =\sum_i\bigl(\max\{\xi_i,0\}-\xi_i\widehat z_i\bigr)
 \le\sum_i|\xi_i|\le n\sigma.
\end{equation}
The bracket holds for every realization of $\xi$, and its width is
additive, not relative. In particular, $(\widehat x,\widehat z)$ is a
feasible solution of the original problem whose cost $\mathrm{UB}$
exceeds $v_0$ by at most $n\sigma$. This makes the role of $\sigma$
explicit: a smaller $\sigma$ narrows the bracket, while the bound
\eqref{eq:explicitwork} grows polynomially with $\sigma^{-1}$.

The certificate is not a cheaper way to approximate $v_0$. For any
$\eta>0$, one call of the oracle of Lemma~\ref{lem:gridoracle} with
$\xi=0$, $I=[n]$, $S=\varnothing$, and accuracy $\eta$ deterministically
returns a support whose exact cost lies in $[v_0,v_0+\eta]$. The
certificate matters when exact perturbed messages or an exact perturbed
solution are computed anyway; it then brackets the original optimum at
no extra cost.
